\documentclass{article}

%----------------------------------------------------------------------------------------
%	PACKAGES AND OTHER DOCUMENT CONFIGURATIONS
%----------------------------------------------------------------------------------------

\usepackage[utf8]{inputenc}
\usepackage{amsmath, amssymb, amsthm, amsfonts}
\usepackage{geometry}
\usepackage{graphicx}
\usepackage{booktabs}
\usepackage{hyperref}
\usepackage{cleveref}
\usepackage{algorithm}
\usepackage{algorithmic}
\usepackage{xcolor}

% Page layout
\geometry{
	a4paper,
	total={170mm,257mm},
	left=20mm,
	top=20mm,
}

% Title and Author
\title{\textbf{Sign-Aware Gated Sparse Autoencoders: \\ Modeling Anticorrelated Features with Bi-Jump-ReLU Activations}}
\author{
    \textbf{Bartosz Wieciech}
}
\date{\today}

% Custom commands for notation
\newcommand{\R}{\mathbb{R}}
\newcommand{\E}{\mathbb{E}}
\newcommand{\gate}{\pi_{\text{gate}}}
\newcommand{\norm}[1]{\left\lVert#1\right\rVert}
\newcommand{\relu}{\text{ReLU}}

\begin{document}

\maketitle

\begin{abstract}
Sparse Autoencoders (SAEs) have emerged as a primary tool for extracting interpretable features from Large Language Models. However, standard SAEs enforce non-negativity on latent activations, compelling the model to learn separate features for diametrically opposed concepts (e.g., ``hot'' vs. ``cold''). This duplication wastes dictionary capacity when features are naturally anticorrelated. We propose the \textit{Sign-Aware Gated SAE}, an extension of the Gated SAE framework that enables a single latent to represent signed evidence along a shared direction. We introduce the \textit{Bi-Jump-ReLU} activation function, which parameterizes a learnable dead zone with independent positive and negative magnitudes. Theoretically, we show that our architecture halves the effective dictionary width required for anticorrelated pairs and becomes strictly parameter-efficient even when such pairs constitute well under $1\%$ of the feature set. We propose a suite of experiments to validate that this approach resolves the "anticorrelation efficiency gap" without sacrificing the no-shrinkage benefits of the gated objective.
\end{abstract}

%----------------------------------------------------------------------------------------
%	1. INTRODUCTION
%----------------------------------------------------------------------------------------

\section{Introduction}

Dictionary learning, particularly via Sparse Autoencoders (SAEs), aims to decompose the dense internal activations of neural networks into sparse, interpretable linear combinations of feature directions \cite{cunningham2023sparse, bricken2023towards}. A standard SAE approximates an activation vector $x \in \R^n$ as $\hat{x} = D f(x) + b$, where $f(x)$ is a sparse vector of coefficients.

A pervasive limitation in current SAE architectures is the non-negativity constraint imposed by the ReLU activation function: $f_i(x) \ge 0$. While valid for many unipolar concepts, this constraint is inductive misalignment for features that represent \textit{anticorrelated} opposites along a single linear subspace---such as sentiment (positive/negative), physical state (hot/cold), or control signals (acceleration/braking). To represent such a subspace, a standard SAE must allocate two distinct latent atoms, $d$ and $-d$, which are mutually exclusive in activation but redundant in geometry. This redundancy inflates the dictionary size required to achieve a given reconstruction fidelity.

Building on the recent success of Gated SAEs \cite{rajamanoharan2024gated}, which decouple feature detection from magnitude estimation to resolve standard $\ell_1$ shrinkage, we introduce the \textit{Sign-Aware Gated SAE}. Our contributions are:

\begin{enumerate}
    \item \textbf{The Bi-Jump-ReLU Activation:} We formalize a signed activation function with a learnable dead zone and asymmetric scaling, allowing a single latent to model "two-sided" evidence.
    \item \textbf{Theoretical Efficiency Analysis:} We prove that our architecture reduces the dictionary width required to model anticorrelated pairs by a factor of 2, with a net parameter savings whenever such pairs exceed roughly $0.2\%$ of the feature set in a 4096-dimensional layer (and an even smaller fraction at larger widths).
    \item \textbf{Protocol Design:} We delineate three experimental protocols---including a reproduction of the geometry from \textit{Toy Models of Superposition} \cite{elhage2022toy}---to rigorously benchmark the approach.
\end{enumerate}

%----------------------------------------------------------------------------------------
%	2. BACKGROUND
%----------------------------------------------------------------------------------------

\section{Background}

\subsection{The Gated SAE and Shrinkage}
Standard SAEs minimize a reconstruction loss combined with an $\ell_1$ penalty on activations:
\[
    \mathcal{L} = \|x - \hat{x}\|_2^2 + \lambda \|f(x)\|_1.
\]
This objective suffers from \textit{shrinkage}: the $\ell_1$ term encourages the model to underestimate the magnitude of active features to lower the penalty, harming fidelity.

Gated SAEs address this by separating the "gate" (detection) from the "magnitude" (estimation). Let $\mu$ denote a fixed centering vector (e.g., a running mean of $x$). The encoder consists of two paths:
\begin{align}
    \pi(x) &= W_{\text{gate}}(x - \mu) + b_{\text{gate}} \\
    f(x) &= \mathbb{I}(\pi(x) > 0) \odot \text{ReLU}\big(W_{\text{mag}}(x - \mu) + b_{\text{mag}}\big).
\end{align}
By applying the $\ell_1$ penalty only to the gate pre-activations $\pi(x)$ and tying the weights (e.g., $W_{\text{mag}} = \mathrm{diag}(\exp r)\, W_{\text{gate}}$ for some per-feature scale vector $r$), Gated SAEs achieve a Pareto improvement in the sparsity-fidelity trade-off.

%----------------------------------------------------------------------------------------
%	3. METHOD
%----------------------------------------------------------------------------------------

\section{Method: Sign-Aware Gated SAE}

We generalize the Gated SAE to handle signed activations while preserving the "no-shrinkage" property. Throughout this section, let $D \in \mathbb{R}^{d_{\text{in}} \times H}$ denote the decoder dictionary with columns $D_{:,i}$, and let $\mu \in \mathbb{R}^{d_{\text{in}}}$ be a fixed centering vector. We use a separate decoder bias $b_{\text{dec}} \in \mathbb{R}^{d_{\text{in}}}$, distinct from the centering.

\subsection{The Bi-Jump-ReLU Activation}
\label{sec:bijumprelu}

We first introduce a shared linear projection for each latent that is geometrically aligned with the decoder:

\begin{equation}
    t_i(x) = D_{:,i}^\top (x - b_{\text{dec}}), \qquad i = 1,\dots,H.
\end{equation}

The gate pre-activation for latent $i$ is then taken to be an affine transformation of this projection:
\begin{equation}
    \pi_i(x) = \alpha_i\, t_i(x) + \beta_i,
\end{equation}
where $\alpha_i, \beta_i \in \mathbb{R}$ are trainable per-latent gate parameters. This makes explicit that detection and magnitude estimation share a direction in input space, following the Gated SAE design.

We introduce two independent threshold parameters $\delta_i^+, \delta_i^- > 0$ (defining a learnable dead zone) and two magnitude scaling factors $r_i^+, r_i^- \in \R$. For notational clarity, we define a polarity variable
\begin{equation}
    s_i(x) =
    \begin{cases}
        +1 & \text{if } \pi_i(x) > \delta_i^+ \\
        -1 & \text{if } \pi_i(x) < -\delta_i^- \\
        0 & \text{otherwise.}
    \end{cases}
\end{equation}

The signed latent activation $a_i(x)$ is then defined as:
\begin{equation}
    \label{eq:bi_jump_relu}
    a_i(x) = 
    \begin{cases} 
        \text{ReLU}\big(e^{r_i^+} t_i(x) + b_{\text{mag}, i}\big) & \text{if } s_i(x) = +1 \\
        -\text{ReLU}\big(-e^{r_i^-} t_i(x) + b_{\text{mag}, i}\big) & \text{if } s_i(x) = -1 \\
        0 & \text{if } s_i(x) = 0,
    \end{cases}
\end{equation}
where $b_{\text{mag},i} \in \mathbb{R}$ is a per-latent offset shared between the positive and negative regimes.

\textbf{Interpretation:} This unit acts as a polarity-sensitive filter. It remains strictly zero within the dead zone $[-\delta_i^-, \delta_i^+]$. If evidence exceeds the positive threshold, it activates with scale $e^{r_i^+}$; if evidence exceeds the negative threshold (in the opposite direction), it activates negatively with scale $e^{r_i^-}$. The decoder column $D_{:,i}$ is shared across both regimes. This effectively collapses two opposite-direction neurons into a single logical unit that carries signed evidence along a common axis.

\subsection{Training Objective}

The decoder reconstructs as
\begin{equation}
    \hat{x} = D a(x) + b_{\text{dec}},
\end{equation}
where $a(x) \in \mathbb{R}^H$ collects the Bi-Jump-ReLU activations from \cref{eq:bi_jump_relu}. To train this structure without shrinkage, we adapt the Gated SAE loss: we place a sparsity-inducing penalty on the gate pre-activations $\pi_i(x)$ and introduce an auxiliary reconstruction task that trains the gate separately from the magnitude parameters.

The total loss per sample is
\begin{equation}
    \mathcal{L}(x) = \underbrace{\|x - D a(x) - b_{\text{dec}}\|_2^2}_{\text{Main Reconstruction}} + \lambda \sum_i \Omega_{\text{gate}}(\pi_i(x)) + \mathcal{L}_{\text{aux}}(x),
\end{equation}
where $\lambda > 0$ is a sparsity strength.

The sparsity term $\Omega_{\text{gate}}$ is a hinge-like, piecewise-linear penalty applied to the gate pre-activations:
\begin{equation}
    \Omega_{\text{gate}}(\pi_i(x)) = \text{ReLU}\big(\pi_i(x) - \delta_i^+\big) + \text{ReLU}\big(-\pi_i(x) - \delta_i^-\big).
\end{equation}
This penalty is exactly zero within the dead zone $[-\delta_i^-, \delta_i^+]$ and grows linearly as $\pi_i(x)$ moves further outside it. Intuitively, it encourages inactive gates to remain inside the dead zone and active gates to exceed their thresholds by the minimum amount necessary; combined with the auxiliary loss, this yields sparse signed activations.

To provide direct supervision to the gate, we introduce an auxiliary reconstruction that uses only the gate output, with a continuous hinge magnitude per side (equal to the amount by which the gate pre-activation exceeds its learned threshold). Let $D^{\text{frozen}}$ denote a copy of the decoder whose parameters are held fixed when optimizing this term (e.g., updated periodically or via an exponential moving average of $D$). The auxiliary loss is
\begin{equation}
    \mathcal{L}_{\text{aux}}(x) = \big\|x - D^{\text{frozen}} \big[ \text{ReLU}(\pi(x) - \delta^+) - \text{ReLU}(-\pi(x) - \delta^-) \big] - b_{\text{dec}}\big\|_2^2,
\end{equation}

\paragraph{Implementation note.}
In our implementation, we stop gradients from the gate sparsity term and the auxiliary reconstruction from flowing into the decoder dictionary. Concretely, the shared projection $t_i(x)$ used to form the gate logits $\pi_i(x)$ is treated as a constant for these terms (stop-gradient), and $D^{\text{frozen}}$ is implemented as the decoder dictionary with gradients detached. This preserves the intended separation: the auxiliary objective trains gate parameters without shaping the decoder geometry.
where $\pi(x)$, $\delta^+$ and $\delta^-$ are applied elementwise across latents.

\textbf{Training dynamics:} The auxiliary reconstruction provides gradients only to the gate-related parameters $(\alpha_i, \beta_i, \delta_i^+, \delta_i^-)$ via $\pi_i(x)$, while $D^{\text{frozen}}$ is not updated through $\mathcal{L}_{\text{aux}}$. The main reconstruction term updates both the decoder $D$ and the magnitude parameters $(r_i^+, r_i^-, b_{\text{mag},i})$ via the signed activations $a_i(x)$. This cleanly separates detection (learned through the auxiliary path and gate penalty) from magnitude estimation (learned through the main reconstruction), mirroring the no-shrinkage argument for Gated SAEs in the non-signed case.

%----------------------------------------------------------------------------------------
%	4. THEORETICAL ANALYSIS
%----------------------------------------------------------------------------------------

\section{Theoretical Analysis}

\subsection{Parameter Efficiency \& Width Reduction}

Let $H$ be the width of a baseline SAE, and suppose that a fraction $p \in [0, 1]$ of the ground-truth features in the data representation appear as perfectly anticorrelated pairs aligned to directions $\pm d$. If there are $N$ total ground-truth features, then $pN$ of them participate in pairs, yielding $(pN)/2$ pairs and $(1-p)N$ unpaired features. A baseline SAE requires two latents for each pair, so it uses $H = N$ latents in total.

In contrast, our Sign-Aware model requires only one latent per anticorrelated pair, plus one latent for each unpaired feature. The effective width required for the Sign-Aware model is therefore
\begin{equation}
    H_{\pm} = (1-p)N + \frac{pN}{2} = N \left( 1 - \frac{p}{2} \right) = H \left( 1 - \frac{p}{2} \right).
\end{equation}
If $p=1$ (all features are paired), the dictionary size halves.

We now analyze the parameter overhead. Let $d_{\text{in}}$ be the input dimensionality. In a standard Gated SAE, each latent carries roughly $d_{\text{in}}$ parameters for the decoder direction plus $O(1)$ scalar parameters for gate biases and scales; we summarize this as $d_{\text{in}} + c$ parameters per latent for some small constant $c$. The Sign-Aware architecture adds 4 scalars per latent (two thresholds, two magnitude scales) compared to the standard Gated SAE, so per-latent parameters become $d_{\text{in}} + c + 4$.

Approximating the total parameter counts as
\begin{align}
    P_{\text{base}} &\approx H (d_{\text{in}} + c), \\
    P_{\text{sign}} &\approx H_{\pm} (d_{\text{in}} + c + 4) = H \left(1 - \frac{p}{2}\right) (d_{\text{in}} + c + 4),
\end{align}
we obtain the condition for net parameter savings:
\begin{align}
    P_{\text{sign}} < P_{\text{base}}
    &\;\Longrightarrow\; \left(1 - \frac{p}{2}\right) (d_{\text{in}} + c + 4) < d_{\text{in}} + c \\
    &\;\Longrightarrow\; 4 - \frac{p}{2}(d_{\text{in}} + c + 4) < 0 \\
    &\;\Longrightarrow\; p > \frac{8}{d_{\text{in}} + c + 4}.
\end{align}
Taking $c \approx 0$ for simplicity yields the approximate threshold
\begin{equation}
    p \gtrsim \frac{8}{d_{\text{in}} + 4}.
\end{equation}
For a model dimension $d_{\text{in}} = 4096$, this gives $p \gtrsim 8/4100 \approx 0.2\%$. Thus, under a simple parameter-counting model, the Sign-Aware SAE is strictly parameter-efficient as soon as anticorrelated pairs constitute even a vanishingly small fraction of the feature set (on the order of $10^{-3}$ for typical transformer layer widths).

\paragraph{Symmetric-magnitude variant.}
An ablation of our architecture ties the per-polarity magnitude scaling parameters ($r_i^+=r_i^-$), while retaining signed gating and asymmetric thresholds. In the parameter-counting model above, this reduces the per-latent scalar overhead by one, i.e., the per-latent count becomes $d_{\text{in}} + c + 3$ instead of $d_{\text{in}} + c + 4$.
Empirically, on Protocol A with the default generator settings (LogNormal positives vs.\ Exponential negatives, $p_+=0.7$), this symmetric-magnitude variant achieves essentially identical reconstruction error and very similar split-regime calibration to the full sign-aware model, indicating that the dominant gain in that setting comes from enabling signed activations along a shared decoder direction rather than from independent per-polarity scaling.

\subsection{Computational Complexity}

\textbf{Inference.} By caching the projection $t(x) = D^\top (x - \mu)$, the forward pass requires only one matrix multiplication for the encoder, identical in structure to the baseline. The subsequent gating and Bi-Jump-ReLU activations operate elementwise in $H_{\pm}$, just as in the standard Gated SAE. Since $H_{\pm} \le H$, the total inference FLOPs scale down by the factor $\left(1 - \frac{p}{2}\right)$ relative to a baseline with width $H$.

\textbf{Training.} The Gated SAE objective requires an additional decoder pass for the auxiliary loss. In a conventional implementation this roughly adds one extra forward/backward pass through the decoder per optimization step. Relative to a baseline SAE with the same target width, this increases decoder-side compute by a factor on the order of $1.5\times$ (the exact factor depends on how encoder and decoder costs balance), and this factor is further reduced by the width multiplier $\left(1 - \frac{p}{2}\right)$ that we gain from exploiting anticorrelated pairs. In other words, training is modestly more expensive per update than a non-gated SAE with the same width, but yields a permanently smaller and faster dictionary at inference time.

%----------------------------------------------------------------------------------------
%	5. PROPOSED EXPERIMENTS (EXTENDED)
%----------------------------------------------------------------------------------------

\section{Proposed Experiments}

We propose a comprehensive evaluation suite designed to isolate the specific contributions of the Sign-Aware Gated SAE: (1) the ability to model asymmetric distributions along a shared axis, (2) the capacity efficiency in anomaly detection tasks, and (3) the geometric consolidation of anticorrelated features.
Unless otherwise stated, we report not only a single operating point, but also sparsity-fidelity trade-offs by sweeping the sparsity strength (e.g., $\lambda$) and plotting reconstruction MSE versus average $L_0$ (active latents per example). This avoids conclusions that depend on a single arbitrary sparsity setting.


\subsection{Protocol A: The "Polarity Dial" (Asymmetry \& Calibration)}
\label{sec:protocol_a}

Protocol A is a controlled synthetic setting where each ground-truth feature is a \emph{signed} scalar coefficient along a
random axis $u_j \in \mathbb{R}^d$. The goal is to isolate whether an SAE can represent and calibrate \emph{both} positive and
negative evidence along the same axis without duplicating decoder atoms. This directly targets the ``anticorrelation efficiency
gap'': non-negative SAEs must allocate two atoms $(+u_j,-u_j)$ to represent one signed axis, while a sign-aware latent can
represent both directions with a single decoder column and a signed activation.

\paragraph{Data generation.}
We construct a synthetic dataset $X \subset \mathbb{R}^d$ consisting of $k$ axes aligned to random directions
$u_j \in \mathbb{R}^d$. Each axis is active with probability $\rho$. If active, its sign $s_j \in \{-1, +1\}$ is sampled from
a Bernoulli distribution with parameter $p_+$, and the magnitude follows different distributions conditioned on the sign:
\begin{align}
    m_j \mid (s_j=+1) &\sim \text{LogNormal}(\mu=0, \sigma=0.5), \\
    m_j \mid (s_j=-1) &\sim \text{Exponential}(\lambda=1.5).
\end{align}
The input is
\begin{equation}
    x = \sum_j s_j m_j u_j + \epsilon,
\end{equation}
with small isotropic noise $\epsilon$.

\paragraph{Baselines and ablations.}
We compare the following model families:
\begin{itemize}
    \item \textbf{Sign-Aware Gated SAE (ours):} Signed latents with two-tailed gating (a learnable dead zone) and the gated
    objective to avoid shrinkage.
    \item \textbf{Sign-Aware, symmetric magnitude (ablation):} The same model, but with tied per-polarity magnitude scales
    ($r_i^+=r_i^-$). This tests whether independent per-polarity scaling is necessary in practice, and provides a more
    parameter-efficient variant.
    \item \textbf{Standard Gated SAE (non-negative, width $H$):} A conventional gated SAE with ReLU latents.
    \item \textbf{Standard Gated SAE (non-negative, width $2H$):} The same gated SAE trained at $2\times$ width. This
    controls for the representational requirement that non-negative models need paired atoms $(+u_j,-u_j)$ to represent a
    signed axis.
    \item \textbf{Baseline ReLU SAE (width $H$):} A standard SAE with ReLU latents.
    \item \textbf{Baseline ReLU SAE (width $2H$):} The same ReLU SAE trained at $2\times$ width.
    \item \textbf{Signed Soft-Threshold SAE (signed-latent baseline):} A signed-latent baseline using a soft-threshold
    nonlinearity (a standard choice in signed sparse coding). This baseline can represent negative coefficients without paired
    atoms, but typically exhibits $\ell_1$ shrinkage and can become dense, separating ``signed latents'' from the gated
    no-shrinkage objective.
\end{itemize}

\paragraph{Axis-aware alignment and pairing.}
To compare learned representations to ground-truth \emph{axes} $u_j$, we evaluate at the axis level rather than treating
$u_j$ and $-u_j$ as distinct targets. We normalize decoder columns to unit norm for alignment, then compute signed cosine
similarities
\[
    s_{ij} = \cos(D_{:,i}, u_j),
\]
and define an axis similarity score $|s_{ij}|$. For sign-aware (signed-latent) models, we perform a one-to-one matching
between latents $i$ and axes $j$ via the Hungarian algorithm on the absolute similarities $|s_{ij}|$, and we record the sign
of $s_{ij}$ for diagnostics.
For non-negative baselines, we additionally allow \emph{two} latents to represent one axis: for each axis $j$, we select the
best positively aligned latent $i^+(j)=\arg\max_i s_{ij}$ and the best negatively aligned latent $i^-(j)=\arg\min_i s_{ij}$
(subject to an alignment threshold). This credits baselines that learn paired atoms $(+u_j,-u_j)$, which is the expected
mechanism by which non-negative models represent signed structure.

\paragraph{Metrics.}
We focus on \textit{Split-Scale Calibration} at the axis level. For each ground-truth axis $j$, let
\begin{equation}
    c_j(x) = s_j m_j
\end{equation}
denote the ground-truth scalar coefficient along $u_j$ in the generative model. We define the model's scalar coefficient
estimate $\hat{c}_j(x)$ by projecting the relevant reconstructed contribution(s) onto $u_j$.
For sign-aware (signed-latent) models with a single matched latent $i(j)$, we use
\begin{equation}
    \hat{c}_j(x) = u_j^\top \big( a_{i(j)}(x)\, D_{:,i(j)} \big).
\end{equation}
For non-negative baselines using paired atoms $(i^+(j), i^-(j))$, we use an axis-aggregated estimate
\begin{equation}
    \hat{c}_j(x) = u_j^\top \Big( a_{i^+(j)}(x)\, D_{:,i^+(j)} + a_{i^-(j)}(x)\, D_{:,i^-(j)} \Big),
\end{equation}
which reduces to $\hat{c}_j(x)\approx a_{i^+(j)}(x)-a_{i^-(j)}(x)$ when $D_{:,i^+(j)}\approx u_j$ and
$D_{:,i^-(j)}\approx -u_j$.
We then compute separate least-squares rescaling factors for the positive and negative halves of the data:
\begin{align}
    \gamma_+ &= \arg\min_{\alpha} \E\big[(\alpha \hat{c}_j(x) - c_j(x))^2 \,\big|\, c_j(x) > 0\big], \\
    \gamma_- &= \arg\min_{\alpha} \E\big[(\alpha \hat{c}_j(x) - c_j(x))^2 \,\big|\, c_j(x) < 0\big].
\end{align}
\begin{itemize}
    \item \textbf{Ideal outcome:} $\gamma_+ \approx 1$ and $\gamma_- \approx 1$.
    \item \textbf{Primary hypothesis:} The Sign-Aware Gated SAE achieves $\gamma_+ \approx 1$ and $\gamma_- \approx 1$ using
    one latent per axis at width $H$. Non-negative models can approach this only if they allocate paired atoms; the $2H$
    baselines test whether remaining gaps at width $H$ are primarily representational (insufficient atoms) rather than
    optimization-related.
    \item \textbf{Secondary question:} Whether independent per-polarity magnitude scaling ($r_i^+ \neq r_i^-$) yields
    additional calibration gains beyond sign-awareness, tested via the symmetric-magnitude ablation.
    \item \textbf{Robustness diagnostics:} Because $\gamma_\pm$ can become ill-conditioned when $\hat{c}_j(x)$ collapses
    toward zero on a regime, we also report per-regime coefficient-space MSE, Pearson correlation between $\hat{c}_j(x)$ and
    $c_j(x)$, and explained variance ($R^2$), each computed separately on $\{c_j(x)>0\}$ and $\{c_j(x)<0\}$.
\end{itemize}

\subsection{Protocol B: Directional Deviation Detection (Resource Constrained)}
\label{sec:protocol_b}

In control systems and operations monitoring, deviations often carry directional semantics (e.g., ``Pressure Too High'' vs.\ ``Pressure Too Low''). Protocol B benchmarks the Sign-Aware Gated SAE as a direction-aware anomaly detector under a strict latent budget constraint, and clarifies when such detection is feasible or fundamentally limited.

\paragraph{Setup.}
We simulate a high-dimensional sensor array $x \in \mathbb{R}^N$ with $N=1024$ channels. We train models under a strict latent budget constraint $H < N$ and evaluate detection at a fixed false positive rate (FPR).

\paragraph{Normal distributions.}
We consider two normal regimes:
\begin{itemize}
    \item \textbf{Isotropic normal:} $x \sim \mathcal{N}(0, I)$.
    \item \textbf{Mildly heteroskedastic normal:} $x_j \sim \mathcal{N}(0, \sigma_j^2)$ with distinct $\sigma_j$ values close to 1. This breaks rotational symmetry while keeping channels approximately comparable, and is intended to make axis semantics learnable without strongly reweighting the reconstruction objective.
\end{itemize}

\paragraph{Anomalies.}
We evaluate two anomaly families:
\begin{enumerate}
    \item \textbf{Coordinate (channel) anomalies:} choose a channel $j$ and shift $x_j \mapsto x_j \pm \Delta$.
    \item \textbf{Feature-direction anomalies:} choose a direction $v_k \in \mathbb{R}^N$ and shift $x \mapsto x \pm \Delta_k v_k$.
\end{enumerate}
For feature-direction anomalies, directions $v_k$ are derived from a reference Sign-Aware model trained on the same normal data. We select the top-$M$ directions by the empirical standard deviation of the reference model's detection signal (gate preactivation) on a clean normal split, which prioritizes directions that produce a strong gate signal. To reduce width-dependent effects, we optionally scale the input perturbation per direction using a fixed-$\Delta \pi$ rule:
\begin{equation}
    \Delta_k = \frac{\Delta}{\alpha_k + \epsilon},
\end{equation}
where $\alpha_k$ is the reference model's per-latent gate scale for the selected latent and $\epsilon$ is a small stabilizer. This targets an approximately constant first-order shift in the reference latent's gate preactivation across selected directions.

\paragraph{Directional scoring and calibration.}
For the Sign-Aware model, let $\pi_i(x)$ denote the gate preactivation of latent $i$. We optionally normalize the per-latent signal on validation normals to improve robustness under heteroskedastic or heavy-tailed normals:
\begin{equation}
    \tilde{\pi}_i(x) = \frac{\pi_i(x) - \mu_i}{\max(\sigma_i, \varepsilon)},
\end{equation}
where $(\mu_i,\sigma_i)$ are fitted on validation normals and $\varepsilon$ is a small floor. We then define exceedances for the positive and negative tails:
\begin{equation}
    z_i^+(x) = \text{ReLU}\big(\tilde{\pi}_i(x) - \tau_i^+\big), \qquad
    z_i^-(x) = \text{ReLU}\big(-\tilde{\pi}_i(x) - \tau_i^-\big),
\end{equation}
where $\tau_i^+$ and $\tau_i^-$ are per-latent thresholds (e.g., the 99th percentile of $\tilde{\pi}_i(x)$ and $-\tilde{\pi}_i(x)$ on validation normals). We pool exceedances across latents to form directional anomaly scores:
\begin{equation}
    \text{Score}_{\text{high}}(x) = \operatorname{Pool}_i\big[z_i^+(x)\big], \qquad
    \text{Score}_{\text{low}}(x)  = \operatorname{Pool}_i\big[z_i^-(x)\big],
\end{equation}
where $\operatorname{Pool}$ is a configurable aggregation (sum, max, or top-$k$ sum). We choose score thresholds $\theta_{\text{high}}$ and $\theta_{\text{low}}$ to achieve a target FPR (1\%) on validation normals, and report recall on held-out anomalies.

\paragraph{Baselines and leakage prevention.}
For baselines with non-negative activations, we use one-sided evidence signals and restrict directional scoring to avoid ``two-sided leakage'' from a single latent. Direction assignment uses decoder geometry (e.g., the sign of the dominant decoder component) and scoring pools exceedances within the assigned direction group.

\paragraph{Metrics.}
We report:
\begin{itemize}
    \item \textbf{Directional Recall @ Fixed FPR:} recall for high and low anomalies at 1\% FPR.
    \item \textbf{Coverage:} the fraction of targets for which both-direction recall exceeds a preset threshold (default 0.8). Targets are channels for coordinate anomalies and directions for feature-direction anomalies.
    \item \textbf{Alignment diagnostics:} decoder channel-purity and the number of unique channels claimed by latents (argmax channel occupancy), used to detect random-rotation behavior and quantify axis alignment.
\end{itemize}

\paragraph{Protocol variants and expectations.}
We treat isotropic coordinate anomalies as a control that highlights non-identifiability under rotational symmetry. We treat symmetry-breaking coordinate anomalies as a stress test of whether channel semantics emerge under a strict latent budget. The primary positive setting is feature-direction anomalies with strength control and localized pooling, which tests whether the Sign-Aware model provides robust two-tailed directional detection on operationally meaningful axes.

\subsection{Protocol C: Geometric Validation on Toy Models}
This protocol replicates the setup of \cite{elhage2022toy} to provide visual and geometric confirmation that the Sign-Aware SAE correctly identifies and consolidates anticorrelated subspaces.

\paragraph{Data Geometry.}
We generate $k$ pairs of features embedded in $\R^2$ (the "unit circle" visualization), with paired directions that are antipodal on the unit circle.
For the qualitative visualizations below, we use the mutually exclusive setting ($\rho=-1$) so that exactly one member of each pair can be active at a time.
We train models with dictionary width $M=k$ (Sign-Aware) and $M=2k$ (Baseline), so that in the ideal case the baseline allocates one latent per feature, while the Sign-Aware model allocates one latent per pair.

\paragraph{Why Sign-Aware Consolidates Antipodal Pairs.}
In the antipodal construction, each ground-truth pair corresponds to a single axis direction $u_j \in \R^2$, with two features at directions $\pm u_j$.
In the mutually exclusive setting ($\rho=-1$), an ideal Sign-Aware SAE can represent the entire pair using a single decoder column aligned with $u_j$:
positive activations correspond to the $+u_j$ feature and negative activations correspond to the $-u_j$ feature.
In contrast, a non-negative SAE cannot use the sign of activations to flip direction, so representing both $\pm u_j$ requires allocating two separate decoder columns, one per feature direction.

\paragraph{Visualization A: The Unit Circle.}
We project the columns of the decoder $D$ onto the 2D data plane and plot them on the unit circle.

\begin{figure}[t]
    \centering
    \includegraphics{protocol_c_unit_circle_baseline_2k_seed0.png}
    \includegraphics{protocol_c_unit_circle_sign_aware_k_seed0.png}
    \caption{Protocol C qualitative geometry: in the mutually exclusive antipodal setting ($\rho=-1$), a non-negative baseline at width $M=2k$ learns one decoder direction per feature, while the Sign-Aware SAE at width $M=k$ learns one direction per pair axis and uses activation sign to distinguish the two antipodal features.}
    \label{fig:protocol_c_unit_circle}
\end{figure}

\begin{itemize}
    \item \textbf{Baseline Expectation:} Two distinct vectors at $180^\circ$ separation (antipodal) for each feature pair.
    \item \textbf{Sign-Aware Expectation:} A single vector per pair. The "negative" feature is represented by the negative activation of the same latent, rather than by an additional decoder vector.
\end{itemize}

\paragraph{Visualization B: The Bi-Jump Histogram.}
We collect activations for a single trained latent over the dataset and plot the histogram of $a_i(x)$.

\begin{figure}[t]
    \centering
    \includegraphics{protocol_c_bijump_hist_sign_aware_k_seed0_logy.png}
    \caption{Bi-jump histogram for a representative Sign-Aware latent in Protocol C. The distribution exhibits two symmetric modes, consistent with a single latent representing an antipodal pair via positive vs negative activations.}
    \label{fig:protocol_c_bijump}
\end{figure}

\begin{itemize}
    \item \textbf{Sign-Aware Expectation:} The histogram should reveal a tri-modal distribution: a large mass exactly at zero (the dead zone $[-\delta_i^-, \delta_i^+]$), a lobe of positive activations, and a lobe of negative activations. This visually confirms the functionality of the Bi-Jump-ReLU activation function defined in \cref{eq:bi_jump_relu}.
\end{itemize}

\paragraph{Quantitative Sweep: Superposition Tolerance.}
To obtain a transition-like quantitative curve, we fix the mutually exclusive antipodal construction ($\rho=-1$) and vary the amount of superposition by changing the number of active pairs per sample (denoted $a$).
As $a$ increases, multiple pairs contribute to each observation, and the model must represent signed pair membership in the presence of superposed features.

\paragraph{Quantitative Metric: Signed Recoverability.}
For each ground-truth pair axis $u_j$, we select the decoder column $D_{:,i}$ with maximum absolute cosine similarity to $u_j$ and define an oriented activation $\tilde a_i(x) = \mathrm{sign}(\langle D_{:,i}, u_j\rangle)\, a_i(x)$.
We then compute the Pearson correlation between $\tilde a_i(x)$ and the signed ground-truth coefficient for that pair, over test samples where the pair is active.
We report the mean of this correlation across pairs as a continuous measure of signed recoverability.

\begin{figure}[t]
    \centering
    \includegraphics{protocol_c_superposition_k16_mean_activation_corr.png}
    \caption{Protocol C quantitative curve (superposition tolerance): signed recoverability decreases sharply as the number of active pairs per sample increases. We plot mean best-candidate activation correlation (mean $\pm$ SEM over seeds).}
    \label{fig:protocol_c_superposition_recoverability}
\end{figure}

\paragraph{Thresholded Metric: Pair Consolidation Rate.}
We also report the \textit{Pair Consolidation Rate}, a stricter, thresholded variant of the above metric: the fraction of ground-truth pairs represented by a single latent that is both axis-aligned and sign-consistent.

We define that a ground-truth pair $(j_1,j_2)$ is \emph{consolidated} by latent $i$ if:
\begin{itemize}
    \item The decoder column $D_{:,i}$ has \emph{absolute} cosine similarity above a threshold $\tau$ (e.g., $\tau = 0.9$) with the pair axis (equivalently, with both antipodal directions), and
    \item The oriented activation $\tilde a_i(x)$ is positively correlated with the signed ground-truth coefficient for the pair, and
    \item The sign of $\tilde a_i(x)$ agrees with the ground-truth sign in both regimes (when $+u_j$ is active vs when $-u_j$ is active).
\end{itemize}
The Pair Consolidation Rate is the fraction of ground-truth pairs for which there exists at least one latent satisfying these conditions.

\begin{itemize}
    \item \textbf{Observation:} Signed recoverability shows a clear transition-like degradation with superposition (Fig.~\ref{fig:protocol_c_superposition_recoverability}). The overcomplete Sign-Aware setting ($M=2k$) degrades faster, consistent with symmetry between one-latent and two-latent solutions.
    \item \textbf{Observation:} The thresholded Pair Consolidation Rate collapses rapidly as superposition increases, indicating that sign-consistent single-latent representations become unreliable under strong interference.
\end{itemize}

\paragraph{Robustness Sweep: Within-Pair Correlation.}
As a robustness check, we also sweep the within-pair correlation $\rho$ from $-1$ (mutually exclusive) to $0$ (independent) at fixed $a=1$ and measure the Pair Consolidation Rate.
This sweep produces a gradual degradation rather than a sharp transition, and is best interpreted as a robustness check rather than a phase transition.

\begin{figure}[t]
    \centering
    \includegraphics{protocol_c_corr_sweep_k8_pair_consolidation.png}
    \caption{Robustness sweep: Pair Consolidation Rate vs within-pair correlation $\rho$ (mean $\pm$ SEM over seeds).}
    \label{fig:protocol_c_corr_sweep}
\end{figure}

%----------------------------------------------------------------------------------------
%	6. RESULTS
%----------------------------------------------------------------------------------------

\section{Results}

\subsection{Protocol A: Asymmetry and Split-Scale Calibration}
\label{sec:results_protocol_a}

\paragraph{Setup.}
We evaluate Protocol A on the synthetic signed-axis dataset described in \S\ref{sec:protocol_a} with
$d=512$ input dimensions and $k=128$ ground-truth axes. Each axis is active with probability $\rho=0.05$ and, when active,
takes a positive sign with probability $p_+=0.7$ and a negative sign otherwise. Conditional on sign, magnitudes are drawn from
a LogNormal$(\mu=0,\sigma=0.5)$ distribution for positives and an Exponential$(\lambda=1.5)$ distribution for negatives, with
additive isotropic noise $\epsilon \sim \mathcal{N}(0, 0.1^2 I)$.
We generate 200k training samples, 20k validation samples, and 20k test samples. All models are trained for 50 epochs with
Adam (learning rate $10^{-4}$, $\beta=(0.9,0.999)$), batch size 1024, and sparsity coefficient $\lambda=10^{-3}$.
We report mean $\pm$ standard deviation over three random seeds (0, 1, 2).

\paragraph{Evaluation protocol.}
We evaluate calibration at the \emph{axis} level using the alignment and aggregation procedure in \S\ref{sec:protocol_a}.
Concretely, we match axes using absolute cosine similarity between decoder columns and ground-truth directions, with a minimum
alignment threshold of $\tau=0.9$. For non-negative baselines, we use the paired-atom aggregation $\hat{c}_j(x)$ when both a
positively aligned and negatively aligned latent are available for a given axis. This gives non-negative models credit for the
intended representation $(+u_j,-u_j)$ and ensures the comparison is not biased against baselines by the evaluation metric.

\paragraph{Metrics.}
We report reconstruction MSE on the test set, the average number of non-zero latents per example ($L_0$), and split-regime
calibration slopes $\gamma_+$ and $\gamma_-$ together with their mean absolute errors. Because $\gamma_\pm$ can become unstable
when a model collapses $\hat{c}_j(x)$ toward zero on one regime, we additionally compute per-regime coefficient-space MSE,
Pearson correlation, and $R^2$; these diagnostics support the same conclusions and are included in the released artifacts.

\begin{table}[t]
    \centering
    \caption{Protocol A results (mean $\pm$ std over three seeds). $H$ denotes the SAE width. $L_0$ denotes the average number
    of non-zero latents per example. $\bar{\gamma}_\pm$ is the mean split-regime calibration slope across axes, and
    $\mathbb{E}[|\gamma_\pm-1|]$ is the mean absolute slope error (lower is better). Non-negative baselines are evaluated with
    axis-level aggregation that credits paired atoms $(+u_j,-u_j)$ when present (\S\ref{sec:protocol_a}).}
    \label{tab:protocol_a_results}
    \small\begin{tabular}{lcccccc}
        \toprule
        Model & MSE$\downarrow$ & $L_0\downarrow$ & $\bar{\gamma}_+$ & $\bar{\gamma}_-$ & $\mathbb{E}[|\gamma_+-1|]\downarrow$ & $\mathbb{E}[|\gamma_--1|]\downarrow$ \\
        \midrule
        ReLU SAE ($H=128$)
            & 0.017987 $\pm$ 0.001168 & 6.62 $\pm$ 0.28
            & 0.780 $\pm$ 0.103 & -1.685 $\pm$ 1.363
            & 0.499 $\pm$ 0.066 & 2.690 $\pm$ 1.371 \\
        ReLU SAE ($H=256$)
            & 0.017138 $\pm$ 0.000286 & 6.88 $\pm$ 0.19
            & 0.913 $\pm$ 0.035 & -1.598 $\pm$ 0.784
            & 0.413 $\pm$ 0.022 & 2.598 $\pm$ 0.784 \\
        Gated SAE ($H=128$)
            & 0.012945 $\pm$ 0.000060 & 9.61 $\pm$ 0.10
            & 0.539 $\pm$ 0.023 & -0.056 $\pm$ 0.011
            & 0.487 $\pm$ 0.023 & 1.056 $\pm$ 0.011 \\
        Gated SAE ($H=256$)
            & \textbf{0.009953 $\pm$ 0.000005} & 9.14 $\pm$ 0.02
            & 1.028 $\pm$ 0.008 & 1.050 $\pm$ 0.007
            & 0.039 $\pm$ 0.010 & 0.060 $\pm$ 0.011 \\
        Signed Soft-Threshold SAE ($H=128$)
            & 0.012038 $\pm$ 0.000241 & 128.00 $\pm$ 0.00
            & 41.387 $\pm$ 1.669 & 24.501 $\pm$ 3.687
            & 43.457 $\pm$ 1.053 & 31.385 $\pm$ 4.389 \\
        Sign-Aware Gated SAE ($H=128$)
            & 0.009963 $\pm$ 0.000004 & 9.15 $\pm$ 0.02
            & 1.036 $\pm$ 0.001 & 1.053 $\pm$ 0.002
            & \textbf{0.036 $\pm$ 0.001} & 0.053 $\pm$ 0.002 \\
        Sign-Aware Gated SAE, symmetric magnitude ($H=128$)
            & 0.009961 $\pm$ 0.000002 & 9.22 $\pm$ 0.10
            & 1.047 $\pm$ 0.014 & 1.045 $\pm$ 0.013
            & 0.047 $\pm$ 0.014 & \textbf{0.045 $\pm$ 0.013} \\
        \bottomrule
    \end{tabular}
\end{table}

\paragraph{The efficiency gap is representational, and sign-awareness closes it.}
At width $H=128$, non-negative models cannot allocate paired atoms $(+u_j,-u_j)$ for all axes, and the deficit shows up as a
systematic failure to calibrate the negative regime. The standard Gated SAE achieves moderate positive-side calibration
($\bar{\gamma}_+ \approx 0.54$) but collapses the negative side ($\bar{\gamma}_- \approx -0.06$), yielding a large negative
regime slope error ($\mathbb{E}[|\gamma_--1|]\approx 1.06$). The ReLU SAE baselines perform substantially worse, reflecting the
combined effect of (i) the same non-negativity constraint and (ii) shrinkage in the standard $\ell_1$ objective.

When we double the width of the non-negative Gated SAE to $H=256$, it can represent each signed axis with two latents and
recovers near-ideal calibration on both regimes ($\bar{\gamma}_+\approx 1.03$, $\bar{\gamma}_-\approx 1.05$).
Crucially, the Sign-Aware Gated SAE reaches essentially the same calibration quality \emph{at width $H=128$} using one latent
per axis. Its reconstruction error (0.009963) is within 0.1\% of the $H=256$ non-negative gated baseline (0.009953), while
using half the dictionary width. This is the intended resolution of the anticorrelation efficiency gap.

\paragraph{Signed activations alone are not sufficient without the gated objective.}
The Signed Soft-Threshold SAE can represent negative coefficients without paired atoms, but it performs poorly in this setting:
it becomes dense ($L_0\approx 128$), exhibits severe miscalibration (very large $\gamma_\pm$ errors), and underperforms the
gated models in reconstruction error. This supports the motivation for combining sign-awareness with the gated no-shrinkage
training objective rather than adopting signed sparse coding activations in an otherwise standard SAE.

\paragraph{Independent $r^+,r^-$ scaling is not a dominant factor at these settings.}
Tying the per-polarity magnitude scales ($r_i^+=r_i^-$) yields virtually identical reconstruction error and very similar
split-regime calibration to the full sign-aware model. In other words, for the default Protocol A generator parameters used
here, the dominant empirical gain comes from enabling signed activations along a shared decoder direction, and we do not
observe a consistent additional benefit from learning separate $r^+$ and $r^-$ values. This suggests the symmetric-magnitude
variant is a reasonable default when prioritizing parameter minimality, while retaining the full model as a strictly more
expressive option for settings with more extreme sign-conditioned asymmetry.

\subsection{Protocol B: Directional Deviation Detection}
\label{sec:results_protocol_b}

\paragraph{Coordinate anomalies on isotropic normals are undetectable under strict compression.}
On isotropic normals $x \sim \mathcal{N}(0, I)$ with single-channel shifts, all evaluated models achieve directional recall close to the calibrated FPR (roughly 1\% to 2\% recall at 1\% FPR) and 0 coverage across widths $H \in \{32,64,128,256,512\}$. Decoder alignment diagnostics are consistent with a random-rotation regime (low channel-purity and argmax-channel occupancy matching random assignment). This matches the expected non-identifiability: when the normal distribution is rotationally symmetric, no coordinate system is preferred, so single-channel shifts do not consistently create large per-latent tail events in the learned basis.

\paragraph{Symmetry-breaking improves alignment but does not rescue coordinate monitoring.}
We next break rotational symmetry using heteroskedastic normals with distinct channel variances, and enable per-latent z-score normalization and max pooling. While alignment diagnostics shift in the expected direction (higher decoder channel-purity and non-random argmax-channel occupancy), directional recall remains near the calibrated FPR and coverage remains 0 across all widths. Empirically, even the best channels detect only a small fraction of anomalies (well below the coverage threshold). This suggests that for independent channels, strict dimensionality reduction $H < N$ can make per-channel monitoring infeasible without additional shared structure (e.g., correlations or low-rank factors) that allows compression without losing channel-level anomaly information.

\paragraph{Feature-direction anomalies yield strong two-tailed detection for the Sign-Aware model.}
We then evaluate feature-direction anomalies aligned to operational axes derived from a reference Sign-Aware model. We select $M=128$ directions by the standard deviation of the reference model's gate preactivation on validation normals (top-$\mathrm{std}(\pi)$), and scale per-direction input perturbations using a fixed-$\Delta \pi$ rule to reduce width-dependent effects. We calibrate thresholds at 1\% FPR and use localized pooling (max or top-$k$).

\begin{table}[t]
    \centering
    \caption{Protocol B results for feature-direction anomalies (mean over three seeds, target FPR 1\%). Directions: $M=128$ selected by top-$\mathrm{std}(\pi)$ from a reference Sign-Aware model; per-direction inputs scaled by a fixed-$\Delta\pi$ rule. Coverage denotes the fraction of directions with both-direction recall $\ge 0.8$.}
    \label{tab:protocol_b_feature_direction}
    \small\begin{tabular}{lcccc}
        \toprule
        Width $H$ & Pooling & Sign-Aware Recall (high) & Sign-Aware Recall (low) & Coverage \\
        \midrule
        128 & max & 0.998 & 0.997 & 1.00 \\
        256 & max & 0.995 & 0.995 & 1.00 \\
        512 & max & 0.992 & 0.990 & 1.00 \\
        \midrule
        512 & top-$k$ ($k=4$) & 0.934 & 0.934 & 1.00 \\
        \bottomrule
    \end{tabular}
\end{table}

Across all widths, the Sign-Aware model achieves near-saturated directional recall and full direction coverage with max pooling (\cref{tab:protocol_b_feature_direction}). In contrast, non-sign-aware baselines remain near chance on this evaluation (directional recall close to the calibrated FPR and 0 coverage), consistent with their one-sided evidence and the difficulty of representing both tails of a direction efficiently under a fixed latent budget.


%----------------------------------------------------------------------------------------
%	References and Appendix would follow
%----------------------------------------------------------------------------------------

\bibliographystyle{alpha}
% \bibliography{references} % Placeholder for BibTeX

\end{document}
